%%%%%%%%%%%%%%%%% PREAMBLE %%%%%%%%%%%%%%%%%%%%%%%%%%%%
\documentclass[letterpaper,10pt,oneside]{article}
\usepackage[utf8]{inputenc}
\usepackage{setspace}
\usepackage{hyperref}
\usepackage{textcomp}
\usepackage{verbatim}
\usepackage{enumitem}
\usepackage{sectsty}
\usepackage{xcolor}
\hypersetup{
	colorlinks,
	linkcolor={red!80!black},
	citecolor={blue!50!black},
	urlcolor={blue!80!black}
}
\sectionfont{\Large}
\subsectionfont{\large}
\usepackage{titlesec}

\renewcommand\labelitemi{--}


\usepackage[left=0.4in, right=0.4in, bottom=0.35in, top=0.35in]{geometry}
\renewcommand{\familydefault}{\sfdefault}
\pagenumbering{gobble}

%\newcommand{\note}[1]{\textcolor{red}{\,(#1)}}
\newcommand{\note}[1]{}

\usepackage{setspace}
\renewcommand{\baselinestretch}{0.94} 

\titlespacing*{\section}{0px}{6px}{4px}
\titlespacing*{\subsection}{0pt}{6px}{8px}
\titlespacing*{\paragraph}{0pt}{6px}{8px}

\input{glyphtounicode}

\pdfgentounicode=1

\begin{document}
{\raggedleft

\begin{minipage}[c]{0.24\textwidth}
	\begin{flushleft}
nil @ mamano . com\\
	\href{http://www.nilmamano.com}{nilmamano.com}
	\end{flushleft}
\end{minipage}\hfill
\begin{minipage}[c]{0.4\textwidth}
	\begin{center}
		\Huge{\textbf{Nil Mamano, PhD}}
	\end{center}
\end{minipage}\hfill
\begin{minipage}[c]{0.24\textwidth}
	\begin{flushright}
	\href{http://www.linkedin.com/in/nilmamano}{linkedin.com/in/nilmamano}\\
	\href{http://www.github.com/nmamano}{github.com/nmamano}
	\end{flushright}
\end{minipage}\\
\vspace{5px}
\hrulefill
}
%\vspace{7px}

%\noindent Computer Science Ph.D. looking to build my career as a software engineer. My research is about algorithms and data structures, specializing in performance. Strong problem-solving skills from solving algorithmic research problems; Effective communication skills from collaborating with 16 co-authors; solid CS fundamentals; Quick learner eager to learn and grow.

\section*{Experience}
\subsection*{Beyond Cracking the Coding Interview}
\vspace{-3px}
\begin{description}
	\item[Author \& Lead Developer]\quad \href{https://www.amazon.com/dp/195570600X}{amazon.com/dp/195570600X}\hfill 2025
	
	Co-wrote the official sequel to Cracking the Coding Interview with Gayle L. McDowell et al. \#1 Best Seller in Data Structures \& Algorithms on Amazon.

	\vspace{2px}
	\noindent Built \href{https://dsatoolkit.com}{dsatoolkit.com}, a companion DS\&A toolkit. Built an automated pipeline to parse, test, and generate articles for all 1,440 code solutions across 5 languages. (Next.js, TypeScript, Python)

	\vspace{2px}
	\noindent Built LLM orchestration pipelines for large-scale codebase changes, e.g., translating 288 book solutions from Python to Go with few-shot prompting and test-driven retry loops, using local models and APIs.

	\vspace{2px}
	\noindent Co-designed an AI interviewer for the book problems (\href{https://bctci.co/ai}{bctci.co/ai}).

\end{description}
\subsection*{Google}
\vspace{-3px}
\begin{description}
	\item[Senior Software Engineer]\hfill Feb 2021 -- Aug 2024

	Worked on Google's software-defined WAN, which carries traffic between all Google data centers. (C++, Go, Python, SQL)

	\vspace{2px}
	\noindent Part of a team that migrated the bandwidth ordering system from a quarterly manual process to an on-demand service with daily resolution, giving internal teams more flexibility and improving WAN utilization.

	\vspace{2px}
	\noindent Built and owned the bandwidth prioritization component. Designed time-series-based algorithms for distributing the available bandwidth among teams when demand exceeds capacity, balancing prioritization, fairness, and continuity constraints.

	\vspace{2px}
	\noindent Part of a cross-functional effort to rapidly provision network capacity for Gemini.

\end{description}

\subsection*{Pathrise}
\vspace{-3px}
\begin{description}
	\item[Tech Interview Consultant]\hfill Apr 2020 -- Jan 2021
	
	Revamped the DS\&A curriculum for coding interviews. Taught algorithms to 100+ students.
\end{description}
\section*{Education}
\begin{description}
\item[PhD + Masters in Computer Science]\quad University of California Irvine, GPA 3.83/4 \hfill Sep 2015 -- Dec 2019

Co-authored 9 peer-reviewed papers on algorithm design, including as main author in tier A conferences like ICALP and ISAAC. The papers describe new algorithmic improvements in graph theory, computational geometry, and computational biology (\href{https://scholar.google.bg/citations?hl=en&user=LIuIigEAAAAJ}{scholar.google.bg/citations?user=LIuIigEAAAAJ}).

\vspace{2px}
\noindent Led a research project from inception to publication: came up with an original problem, engaged 3 colleagues to work on it, and collaborated with them to solve it and write a paper. We invented an algorithm for the knight's tour problem (cited by Knuth).

\vspace{2px}
\noindent Led 100+ sessions and guest lectures with 50+ students. Led a 120-student study on the effect of immediate automated feedback on learning outcomes.

\item[B.E. in Computer Science]\quad Polytechnic University of Catalonia, GPA 3.8/4 (99th percentile) \hfill Sep 2011 -- Jul 2015

Created SANA, a C++ biological network alignment tool that uses simulated annealing. 100+ citations; actively maintained for 10+ years by 50+ collaborators, with 2,000+ commits (\href{https://github.com/nmamano/SANA}{github.com/nmamano/SANA}).
\end{description}

\section*{Projects}
\begin{description}
	\item[Wall Game]\quad \href{https://wallgame.io}{wallgame.io} \hfill 2025

	A full-stack online board game supporting AIs as first-class citizens via an AI integration protocol/client. (React, TypeScript, Tailwind, shadcn, PostgreSQL, MongoDB, WebSocket)

	\vspace{2px}
	\noindent Trained multiple AlphaZero-style models for different board sizes and rule sets via self-play on consumer GPUs. (C++, PyTorch, TensorRT, ONNX, WebAssembly)

	\item[Technical Blog]\quad \href{https://nilmamano.com/blog}{nilmamano.com/blog} \hfill 2025

	30+ articles on algorithms, data structures, AI, CS research, and software engineering.

\begin{comment}
	 
	\item[SANA 2.0]\quad  \href{https://github.com/nmamano/SANA}{github.com/nmamano/SANA} \hfill Apr 2020 -- Jun 2020
	
	C++ $|$ Extended the SANA software (mentioned above) to enable alignment-finding between viruses. Refactored 20k+ lines of the codebase and added new features, allowing academic collaborators to run scientific experiments comparing SARS-CoV-2 to other viruses (currently underway).

	\item[ttl-cache]\quad  \href{https://github.com/nmamano/ttlcache}{github.com/nmamano/ttlcache}\hfill  Mar 2020
	
	C++ (templates and C++17 features) $|$ An efficient in-memory key-value cache that supports timeouts. It implements Redis' algorithm to remove expired entries and has an LRU eviction mechanism.
	
	\item[Graph Nearest-Neighbor]\quad  \href{https://github.com/nmamano/NearestNeighborInGraphs}{github.com/nmamano/NearestNeighborInGraphs} \hfill 2018
	 
	 Java $|$ Designed and implemented data structures for matching drivers and passengers on a road network dynamically. Leveraged the topology of road networks to optimize the data structures and scale them to statewide networks with 100k+ edges.  
	\item[Stable Matching Voronoi Diagram]\quad  \href{https://github.com/nmamano/StableMatchingVoronoiDiagram}{github.com/nmamano/StableMatchingVoronoiDiagram} \hfill 2018
	
	C++, OpenGL, Javascript $|$ Designed and implemented clustering algorithms for partitioning data into balanced clusters, and researched their application to gerrymandering prevention: they can partition a geographic region into compact and equal-population districts. Built a desktop application to visualize and benchmark the algorithms.
\end{comment}	
\begin{comment}
	\item[Chespel]\quad  \href{https://github.com/nmamano/Chespel}{github.com/nmamano/Chespel}\hfill 2014
	
	Java $|$ Designed a high-level programming language to specify the position evaluation rules of a chess engine. Built a transpiler that converts the code into C++ and integrates it with the search algorithm from the Faile chess engine, allowing users to build their own chess engines and experiment with the evaluation rules without needing to understand the search algorithm.
\end{comment}	
\end{description}




\end{document}
\begin{comment}
\subsection*{Pathrise}

\paragraph*{Technical Advisor} \hfill Apr 2020 -- Present

\vspace{-1px}
\noindent ((Still writing this.))

%\begin{itemize}[leftmargin=15px]
%	\item 2nd highest GPA in a cohort of 400+ students.
%\end{itemize}



%\item Improved the runtime of the multi-fragment algorithm for Euclidean TSP from $O(n^2)$ to $O(n\log n)$.
%	\item Designed a parallel algorithm for DNA sequencing.
%\item Proved that a class of Greedy algorithms can be computed in a distributed system.
%\item Implemented an algorithm for embedding hidden ``watermarks'' in large graphs by hiding information in the structure of the graphs, allowing companies to share data encoded as graphs such as social networks without the risk of it being leaked. [Python]


%\section*{Selected publications}

Selected publications:
\begin{itemize}[leftmargin=15px]
%\item ``Euclidean TSP, Motorcycle Graphs, and Other New Applications of Nearest-Neighbor Chains'' N. Mamano, A. Efrat, D. Eppstein, D. Frishberg, M.T. Goodrich, S. Kobourov, P. Matias, V. Polishchuk, International Symposium on Algorithms and Computation 2019
\item ``Stable-Matching Voronoi Diagrams: Combinatorial Complexity and Algorithms'' G. Barequet, D. Eppstein, M. T. Goodrich, and N. Mamano, International Colloquium on Automata, Languages and Programming 2018
\item ``Reactive Proximity Data Structures for Graphs'' D. Eppstein, M.T. Goodrich, and N. Mamano, Latin American Theoretical Informatics Symposium 2018
\item ``SANA: Simulated Annealing far outperforms many other search algorithms for biological network alignment'' N. Mamano and W. Hayes, Bioinformatics: Oxford Journals 2017
\end{itemize}
\end{comment}
\begin{comment}
\begin{itemize}
%	\item Designed a parallel algorithm for DNA sequencing.
%\item Proved that a class of Greedy algorithms can be computed in a distributed system.
%\item Implemented an algorithm for embedding hidden ``watermarks'' in large graphs by hiding information in the structure of the graphs, allowing companies to share data encoded as graphs such as social networks without the risk of it being leaked. [Python]

\end{itemize}
\end{comment}

%\section*{Awards and Honors}
%\begin{itemize}
%	\item Balsells Graduate Fellowship\hfill 2015 -- 2016
%	\item Balsells Fellowship---Undergraduate Mobility Program\hfill 2014
%	\item Formula Santander International Scholarship \hfill 2014
%\end{itemize}
